% =====================================================================
\section{Grundgerüst und Text}
% =====================================================================

\begin{frame}[fragile]{Das minimale Dokument}
  \begin{columns}[c]
    \begin{column}{0.56\textwidth}
\begin{code}
\documentclass{article}   % Dokumenttyp

% Präambel: Einstellungen und Pakete
\usepackage[ngerman]{babel}
\usepackage{amsmath}

\title{Mein erstes Dokument}
\author{Max Mustermann}

\begin{document}          % Inhalt
\maketitle
Hallo Welt!
\end{document}
\end{code}
    \end{column}
    \begin{column}{0.4\textwidth}
      \begin{itemize}
        \item \textbf{Präambel:} alles vor \texttt{\textbackslash begin\{document\}}
        \item Befehle beginnen mit \texttt{\textbackslash}
        \item Pflichtargumente in \texttt{\{\}}, optionale in \texttt{[\,]}
        \item \texttt{\%} leitet einen Kommentar ein
      \end{itemize}
    \end{column}
  \end{columns}
\end{frame}

\begin{frame}{Wichtige Pakete für den Anfang}
  \centering
  \rowcolors{2}{papier}{hell}
  \begin{tabular}{@{}p{4.6cm}p{6.8cm}@{}}
    \toprule
    \textbf{Paket} & \textbf{Wofür?}\\
    \midrule
    \texttt{babel} mit \texttt{ngerman} & Deutsche Silbentrennung und Bezeichnungen\\
    \texttt{amsmath}, \texttt{amssymb} & Erweiterter Mathesatz und Symbole\\
    \texttt{siunitx} & Zahlen und Einheiten korrekt setzen\\
    \texttt{graphicx} & Bilder einbinden\\
    \texttt{booktabs} & Schöne Tabellen\\
    \texttt{hyperref} & Klickbare Verweise im PDF\\
    \bottomrule
  \end{tabular}
\end{frame}

\begin{frame}[fragile]{Struktur, Listen und Verweise}
  \begin{columns}[c]
    \begin{column}{0.53\textwidth}
\begin{code}
\section{Einleitung}
\label{sec:einleitung}

Text mit \textbf{fett} und \emph{betont}.

\begin{itemize}
  \item erster Punkt
  \item zweiter Punkt
\end{itemize}

\subsection{Messung}
Siehe Abschnitt~\ref{sec:einleitung}.
\end{code}
    \end{column}
    \begin{column}{0.43\textwidth}
      \begin{itemize}
        \item Abschnitte werden automatisch nummeriert
        \item \texttt{itemize} für Stichpunkte, \texttt{enumerate} für Nummern
        \item \texttt{\textbackslash label} und \texttt{\textbackslash ref}: Verweise passen sich selbst an
        \item Leerzeile beginnt einen neuen Absatz
        \item \texttt{\textasciitilde} ist ein geschütztes Leerzeichen
      \end{itemize}
    \end{column}
  \end{columns}
\end{frame}

\begin{frame}[fragile]{Fehlermeldungen lesen}
\begin{code}[listing options={style=log}, colback=akzent!10!papier, colframe=akzent!10!papier]
! Missing $ inserted.
<inserted text>
                $
l.12 Die Energie ist E = mc^
                            2.
\end{code}
  \vspace{0.3em}
  \begin{columns}[T]
    \begin{column}{0.47\textwidth}
      \begin{itemize}
        \item \texttt{l.12}: Fehler in Zeile 12
        \item Hier steht \texttt{\^{}} außerhalb des Mathemodus, es fehlt \texttt{\$\,\dots\$}
      \end{itemize}
    \end{column}
    \begin{column}{0.47\textwidth}
      \begin{merke}
        Immer den \textbf{ersten} Fehler beheben, dann neu kompilieren.
      \end{merke}
    \end{column}
  \end{columns}
\end{frame}
